\chapter{牛顿宇宙学与弗里德曼方程}

上次我们讨论了空间的度规。接下来的问题是动力学：标度因子如何随时间变化？首先讲一点几何，然后是牛顿宇宙学。弗里德曼方程暂时还不会从广义相对论推导出来，而是将从牛顿方程推导出来。

\section{均匀性、各向同性与含时度规}

均匀性和各向同性是基本假设。如果宇宙是均匀且各向同性的，并且随时间变化，那么哈勃定律总是成立。在任意固定时刻，退行速度与距离成正比，
\begin{equation}
v = H(t)\,d.
\end{equation}
这种正比关系是线性的。恒定的不是时间，而是空间：在给定瞬间，同一个比例因子适用于各处。

\begin{classroomqa}
\audiencequestion{这是说线性正比吗？}\lecturetimestamp{00:02:06}
\lecturerresponse{是的。在固定时刻，退行速度与距离成线性正比。比例因子就是哈勃常数。一般来说它依赖于时间，但它在空间中不随位置变化。}
\end{classroomqa}

这就是为什么``哈勃常数''这个说法是危险的。它相对于空间位置是常数，但相对于时间不是常数。

接下来是一个自然的问题。如果空间在膨胀，为什么米尺不也跟着膨胀？答案是，把米尺结合在一起的力，远远强于原子尺度上与宇宙学膨胀相关的任何效应。束缚物体保持束缚。膨胀的是足够不相互束缚的物体之间的距离，例如相距足够远、彼此不能强烈地束缚在一起的星系。

\begin{classroomqa}
\audiencequestion{测量杆也在膨胀，这会不会有问题？}\lecturetimestamp{00:04:00}
\lecturerresponse{不会。测量杆并不膨胀。膨胀的是坐标。把米尺结合在一起的力远远强于膨胀引起的效应，所以普通束缚物体不会随背景膨胀。相距足够远的星系确实会彼此分离，但米尺不会。}
\end{classroomqa}

甚至可以设想原子之间存在一种微小的额外排斥效应。这仍然不会把一根杆撕开。它只会稍微移动平衡间距，就像一点额外的拉力会改变一根刚性弹簧的平衡位置一样。

\begin{classroomqa}
\audiencequestion{不过，你不能把它建模为点之间的力，对吗？因为否则它就会是一种加速度。}\lecturetimestamp{00:05:29}
\lecturerresponse{可以设想原子之间存在一种微小的额外力，但如果它相对于结合力非常微小，它就只会稍微移动平衡位置。膨胀本身并不会产生一种普通的力律。某些类型的膨胀可以用力的方式来建模，也就是加速膨胀，但那是后面才讨论的问题。}
\end{classroomqa}

所以有用的坐标是嵌入在膨胀物质分布中的坐标。没有额外本动运动的星系保持相同的坐标值。随时间变化的是坐标间距与以米测量的物理距离之间的关系。这种关系由标度因子\(a(t)\)编码。

对于平坦空间，空间度规可以写成
\begin{equation}
ds^2 = a(t)^2\left(dx^2+dy^2+dz^2\right).
\end{equation}
这里\(a\)是相隔一个单位坐标距离的相邻点之间的实际距离，以米为单位。如果宇宙膨胀，\(a\)就变得依赖于时间。

要过渡到时空，从狭义相对论形式出发
\begin{equation}
d\tau^2 = dt^2 - ds^2,
\end{equation}
并代入空间度规。这给出
\begin{equation}
d\tau^2 = dt^2 - a(t)^2\left(dx^2+dy^2+dz^2\right).
\end{equation}
这些系数不依赖于空间中的位置。它们只依赖于时间。这就是用度规语言表述的均匀性。对角度规具有如下形式
\begin{equation}
g_{\mu\nu}=\mathrm{diag}\!\left(1,-a(t)^2,-a(t)^2,-a(t)^2\right).
\end{equation}
\editorialnote{这里写出的完整对角矩阵遵循00:14:33附近口述的度规张量描述。黑板画面清楚地显示了线元，但没有显示完全清晰可读的矩阵。}

没有非对角项。没有\(dx\,dy\)，没有\(dx\,dt\)，也没有其他混合项。这个度规几乎和狭义相对论时空的度规一样简单。特殊之处在于，空间度量与时间之间的关系本身依赖于时间。

\begin{figure}
\centering
\includegraphics[width=0.62\textwidth]{lecture_02_figure_02.png}
\caption{共动网格草图以及含时标度因子的时空间隔。}
\end{figure}

这幅图说明了要点。Alice 和 Bob 在共动坐标中始终相距一个单位，但随着时间推移，这一个坐标单位代表的米数越来越多。如果在过去的某个时刻 \(a(t)\) 为零，那么每一个物理间距都会是零。一切都会被压叠在其他一切之上。这就是为什么 \(a\to 0\) 对应无限密度。如果 \(a\) 增大，宇宙就变得越来越稀薄。

度规的这种形式并非在每个参考系中都成立。它是在随星系运动的参考系中的形式。

\begin{classroomqa}
\audiencequestion{在那个方程中，小的 \(t\)，\(a(t)\)，它的 \(t\) 是什么？}\lecturetimestamp{00:19:19}
\lecturerresponse{它是由随星系运动的时钟所测得的时间。想象在同一家宇宙时钟工厂制造、并由星系携带的相同时钟。在这个参考系中，时间变量正是那些时钟所测量的量。沿星系世界线，固有时与这个坐标时间相同。}
\end{classroomqa}

确实，沿共动星系世界线，
\begin{equation}
dx=dy=dz=0,
\end{equation}
所以
\begin{equation}
d\tau^2 = dt^2,
\end{equation}
因此
\begin{equation}
dt=d\tau.
\end{equation}

\begin{classroomqa}
\audiencequestion{是否存在使度规的这种形式保持不变的变换？}\lecturetimestamp{00:23:08}
\lecturerresponse{旋转可以，因为 \(dx^2+dy^2+dz^2\) 在旋转下不变。但洛伦兹助推不保持这种形式。物质本身挑出了随星系运动的参考系，而在其他某个参考系中，世界看起来不会是各向同性的。}
\end{classroomqa}

\section{引力场中的抛体}

在把牛顿定律应用于整个宇宙之前，先退回到一个更简单的问题会有帮助：一个重质量 \(M\)、一个远小于它的质量 \(m\)，以及一个初始的向外冲量。在冲量之后，小质量只是在引力场中运动，并按照牛顿定律所要求的那样行动。

对于宇宙学讨论而言，重要的不是完整的轨道问题，而是能量守恒。动能是
\begin{equation}
\frac{1}{2}mv^2.
\end{equation}

\begin{figure}
\centering
\includegraphics[width=0.42\textwidth]{lecture_02_figure_03.png}
\caption{向外运动的试验质量与动能项 \(\frac{1}{2}mv^2\)。}
\end{figure}

引力势能包含两质量的乘积、牛顿常数，以及分母中距离的一次幂。板书上这些因子以 \(mMG/D\) 的顺序出现，所以这里最好保持那个顺序：
\begin{equation}
-\frac{mMG}{D}.
\end{equation}
因此总能量是
\begin{equation}
\frac{1}{2}mv^2-\frac{mMG}{D}.
\end{equation}
能量守恒只是说，对于所讨论的这个特定粒子，这个量不随时间改变。不同的粒子可以有不同的能量。

\begin{figure}
\centering
\includegraphics[width=0.62\textwidth]{lecture_02_figure_04.png}
\caption{径向运动的试验质量的总牛顿能量：动能项减去引力势能。}
\end{figure}

负号很重要。引力势能是负的，因为引力是吸引的。把粒子拉近会产生更大的负势能。把它移得更远则使势能负得更少。

\begin{classroomqa}
\audiencequestion{势能随距离增加而不是减少？}\lecturetimestamp{00:34:13}
\lecturerresponse{是的。忘掉常数，看看 \(1/D\) 和 \(-1/D\) 的图像。力推向势能减小的方向。因此吸引的力对应倒过来的曲线 \(-1/D\)，所以势能随距离增加，并且当粒子向内运动时变得更负。}
\end{classroomqa}

\begin{figure}
\centering
\includegraphics[width=0.46\textwidth]{lecture_02_figure_05.png}
\caption{关于引力势能符号的 \(1/D\) 与 \(-1/D\) 的定性比较。}
\end{figure}

\section{逃逸速度}

现在对可能性进行分类。能量可以是正的、负的或零。只有这三种情况。

\begin{itemize}
\item 如果能量为负，粒子是束缚的。以过小的速度把它向上抛出，它会上升、减速、掉头，然后落回来。
\item 如果能量为正，粒子会逃逸。引力起初使它减速，但不足以使它停下，最终它以非零的渐近向外速度离去。
\item 如果能量恰好为零，粒子处于逃逸速度。那是逃逸与不逃逸之间的临界情况。
\end{itemize}

对于足够小的初始速度，势能项从一开始就占主导，能量为负。

\begin{figure}
\centering
\includegraphics[width=0.62\textwidth]{lecture_02_figure_06.png}
\caption{在小的向外速率下，负的势能项在总能量中占主导。}
\end{figure}

在零能量情况下，距离不断增大，但速率不断减小。粒子不可能在有限时间内静止下来。如果它静止下来，就会落回去。所以距离是无界的，而速度趋于零。

这就是对宇宙学至关重要的直接类比。宇宙可以以同样的方式表现。如果初始向外的哈勃流足够大，那么一切都会逃离其他一切。如果它太小，引力就会获胜，并发生一次巨大的坍缩。还存在临界情形，它恰好类似于逃逸速度。

\section{牛顿定理}

宇宙学上的简化之所以成立，是因为平方反比定律的牛顿定理。

如果物质围绕某个中心呈球对称分布，那么对于分布外的一个粒子，引力恰好与全部质量都集中在中心时相同。但该定理比这更强。如果粒子位于分布内部，那么只有半径更小处的质量有贡献。半径更大处的质量不给出净力。

该定理最简洁的版本是薄壳。一个完美的球壳对其内部任意位置的粒子都不施加引力，即使不在中心处也是如此。乍看之下这令人惊讶，但对于\(1/r^2\)定律而言这是真的。

\begin{classroomqa}
\audiencequestion{在你那里给出的例子中，我们是否应该有点惊讶：当粒子穿过边界时，力似乎是间断的？}\lecturetimestamp{00:48:19}
\lecturerresponse{是的。对于理想的无限薄壳，这令人惊讶，而且这是\(1/r^2\)力定律所特有的。真实的壳有厚度，因此变化会被抹平，但对于理想壳的尖锐陈述正是牛顿定理的非凡性质。}
\end{classroomqa}

一幅图像随后使该定理显得合理。把一个点质量替换为一个点水源。在三维中，向外的流量按\(1/r^2\)减小。把这样的源连续地放在一个球壳周围各处。在壳外，流动看起来恰好就像一切都集中在中心一样。在壳内不可能有净流动，因为中心没有源。正是同样的结构构成了引力定理的基础。

这里的要点不是重新推导该定理，而是使用它。

\section{作为单体问题的均匀宇宙}

现在回到宇宙学。取世界的一小块区域，其中充满均匀的星系气体。如果愿意，可以把观察者自己放在中心。这只是一种方便。任何人都可以自称位于中心，而必须得到同一个方程。

选择一个位于\(x\)轴上、共动坐标为\(x\)的星系。由于宇宙是各向同性的，轴的选择无关紧要。画一个以\(x=0\)为中心并通过该星系的球。根据牛顿定理，作用在该星系上的力恰好与球内全部质量都集中在中心时相同，而球外的一切都不贡献。

这个球随星系一起膨胀。它始终包含同样的星系，因此包含同样的总质量，尽管其物理半径随时间变化。

因此，该星系的运动恰好遵循与抛射体相同类型的能量方程：
\begin{equation}
\frac{1}{2}mv^2-\frac{mMG}{D}=\text{constant}.
\end{equation}
这里\(M\)现在指的是膨胀球内的质量，而不是热身图像中的中心质量。随后，通过乘以\(2\)并除以\(m\)来简化代数。由于一个常数除以\(m\)仍然只是一个常数，因此可以写出
\begin{equation}
v^2-\frac{2MG}{D}=\text{constant}.
\end{equation}
同样，``常数''意味着对该特定星系而言不随时间变化。

现在把它与标度因子联系起来。到该星系的物理距离是
\begin{equation}
D=ax,
\end{equation}
因为\(x\)是固定的共动坐标。它的速度是
\begin{equation}
v=\dot D=\dot a\,x.
\end{equation}

球内的质量是密度乘以体积。如果\(\rho\)是质量密度，那么
\begin{equation}
M=\frac{4\pi}{3}\rho D^3=\frac{4\pi}{3}\rho a^3x^3.
\end{equation}
代入能量方程得到
\begin{equation}
\dot a^2x^2-\frac{2G}{ax}\left(\frac{4\pi}{3}\rho a^3x^3\right)=\text{constant},
\end{equation}
或
\begin{equation}
\dot a^2x^2-\frac{8\pi G}{3}\rho a^2x^2=\text{constant}.
\end{equation}

左侧两项都与\(x^2\)成正比。那里其他的一切在各处都相同。因此右侧也必须与\(x^2\)成正比。这与把它称为常数并不矛盾。它意味着对每个固定的\(x\)不随时间变化，但它仍可能从一个星系到另一个星系随\(x^2\)缩放。这是合理的。更遥远的星系运动得更快。

所以除以公因子\(x^2\)。剩下的常数现在既与时间无关，也与空间无关。按历史上标准的写法把它写成\(-k\)：
\begin{equation}
\dot a^2-\frac{8\pi G}{3}\rho a^2=-k.
\end{equation}
除以\(a^2\)后，
\begin{equation}
\left(\frac{\dot a}{a}\right)^2=\frac{8\pi G}{3}\rho-\frac{k}{a^2}.
\end{equation}

这就是著名的弗里德曼-罗伯逊-沃克方程，不过讲义也指出它实际上就是弗里德曼方程。符号约定是历史性的。写出\(-k\)并不意味着\(k\)必须为正。\(k\)可以为正、为负或为零。在光速取为\(1\)的单位制中，\(k\)是无量纲的。

\section{密度与物质主导方程}

接下来的问题立刻出现：\(\rho\)是什么？它不是一个常数。随着宇宙膨胀，密度会减小。

为了表达这种依赖关系，在坐标空间中取一个单位立方体：在\(x\)、\(y\)和\(z\)方向上边长均为一个单位。同样的星系始终位于该立方体内部，因为该立方体随坐标一起膨胀。因此，该坐标立方体内的质量是固定的。

\editorialnote{讲义在几个阶段重复使用了符号\(M\)。这里使用\(M_{\Box}\)只是为了将单位坐标立方体内的质量与之前的中心质量和球体内包含的质量区分开来。}

该立方体的物理边长为\(a\)，因此其物理体积为\(a^3\)。因此
\begin{equation}
\rho=\frac{M_{\Box}}{a^3}.
\end{equation}
这就是物质主导情形：右端只涉及普通非相对论性物质的密度。

\begin{classroomqa}
\audiencequestion{根据该方程，哈勃常数是\(\rho\)的函数。因此它不是一个常数。}\lecturetimestamp{01:19:21}
\lecturerresponse{正是如此。哈勃常数不是一个常数。它依赖于密度，而密度随时间变化。讲义甚至开玩笑说，期末考试可以只出一道题：哈勃常数是常数吗？}
\end{classroomqa}

将\(\rho=M_{\Box}/a^3\)代入弗里德曼方程得到
\begin{equation}
\left(\frac{\dot a}{a}\right)^2=\frac{8\pi G M_{\Box}}{3a^3}-\frac{k}{a^2}.
\end{equation}

此时任意选择已经清楚了。需要单位坐标立方体内的质量和\(k\)的值。\(k\)的三种可能符号对应于抛射体问题中的三种能量情形。

\begin{itemize}
\item 如果\(-k>0\)，则能量实际上高于逃逸速度，宇宙将永远膨胀。
\item 如果\(-k<0\)，则引力最终获胜，宇宙将重新坍缩。
\item 如果\(k=0\)，则宇宙恰好处于逃逸速度。
\end{itemize}

\begin{classroomqa}
\audiencequestion{如果\(k=0\)呢？那么两个星系之间的速度趋近于零吗？}\lecturetimestamp{01:21:12}
\lecturerresponse{是的。在临界情形中，速度趋近于零，但距离仍然无界。星系继续分离，只是越来越慢。如果它们真的在某个有限时刻静止下来，它们就会落回去。}
\end{classroomqa}

\section{临界情形\(k=0\)}

现在选择最简单的情形，\(k=0\)。这正是逃逸速度情形，它之所以更简单只是因为数学更简单。方程变为
\begin{equation}
\left(\frac{\dot a}{a}\right)^2=\frac{8\pi G M_{\Box}}{3a^3}.
\end{equation}

右端始终为正且永不为零。因此\(\dot a/a\)不可能穿过零。在这种情形下宇宙不可能掉头。它要么始终膨胀，要么始终收缩。取膨胀分支。

现在猜测一个幂律，
\begin{equation}
a(t)=ct^p.
\end{equation}
于是
\begin{equation}
\dot a = cp\,t^{p-1},
\end{equation}
所以
\begin{equation}
\frac{\dot a}{a}=\frac{p}{t}.
\end{equation}
代入微分方程得到
\begin{equation}
\frac{p^2}{t^2}=\frac{8\pi G M_{\Box}}{3c^3t^{3p}}.
\end{equation}
要使\(t\)的幂次匹配，必须有
\begin{equation}
3p=2.
\end{equation}
所以
\begin{equation}
p=\frac{2}{3}.
\end{equation}
于是常数必须满足
\begin{equation}
\frac{4}{9}=\frac{8\pi G M_{\Box}}{3c^3},
\end{equation}
因此
\begin{equation}
c^3=6\pi G M_{\Box}.
\end{equation}
所以
\begin{equation}
a(t)=\left(6\pi G M_{\Box}\right)^{1/3}t^{2/3}.
\end{equation}

仍然可以有一个附加的时间平移。它通过将\(t\)替换为\(t-t_0\)而被吸收。一旦选择时间原点使得在\(t=0\)时\(a=0\)，那个剩余的常数就消失了。

因此宇宙按时间的三分之二次方膨胀，
\begin{equation}
a(t)\propto t^{2/3}.
\end{equation}
它并不随时间线性膨胀。线性增长将对应于无加速运动。在这里膨胀在减慢。

如果把公式向前推到\(t=0\)，标度因子趋于零，密度发散。说\(a=0\)是极端的。我们不知道当物质变得过于致密时会发生什么。真正的要点在于，在某个早期阶段，一切都如此拥挤，以至于原子、原子核和夸克全都挤在一个极其致密的状态中。

\section{哈勃常数与宇宙年龄}

在临界模型中，
\begin{equation}
H(t)=\frac{\dot a}{a}=\frac{2}{3t}.
\end{equation}
所以哈勃常数近似等于宇宙年龄的倒数，
\begin{equation}
t=\frac{2}{3H}.
\end{equation}
这是该计算的主要收获之一。测量今天的哈勃常数，就能得到宇宙年龄的估计，至少可以回溯到牛顿模型不再适用的那个时刻。

所引用的年龄大致为\(13\)到\(14\)亿年。文中还说，\(k=0\)在观测上非常现实：在文中引用的精度水平上，尚未知道与\(k=0\)有显著偏离。

\begin{classroomqa}
\audiencequestion{我只是有点搞不清楚\(M\)和\(a\)是什么。}\lecturetimestamp{01:39:16}
\lecturerresponse{两者都包含来自网格选择的任意性。改变基准盒，其中的物质量就会改变，标度因子的数值也会改变。物理的是密度\(M/a^3\)，或者等价地，诸如\(a/M^{1/3}\)这样不依赖于网格如何选取的组合。在\(k=0\)解中，乘在\(t^{2/3}\)上的未知常数在\(\dot a/a\)中被消去，因此年龄关系\(t=2/3H\)不依赖于基准盒的质量。}
\end{classroomqa}

这种消去在解中是直接可见的。如果
\begin{equation}
a(t)=Ct^{2/3},
\end{equation}
那么
\begin{equation}
\dot a=\frac{2}{3}Ct^{-1/3},
\end{equation}
因此
\begin{equation}
\frac{\dot a}{a}=\frac{2}{3t}.
\end{equation}
常数\(C\)完全消失。

\section{为什么牛顿推导有效}

接近末尾时，问题变成：为什么牛顿推导会与相对论推导如此一致？

\begin{classroomqa}
\audiencequestion{这一切似乎都依赖于一个开放宇宙的假设。}\lecturetimestamp{01:45:44}
\lecturerresponse{在这个阶段，推导是牛顿式的，而牛顿物理意味着平直空间。\(k\)在空间几何方面的解释被推迟到下一讲。在那里，\(k=0\)将对应于平直空间，而其他符号将与弯曲的空间几何相联系。}
\end{classroomqa}

\begin{classroomqa}
\audiencequestion{很久以前，米尺的长度是否依赖于\(a\)？}\lecturetimestamp{01:47:53}
\lecturerresponse{足够早的时候，物质会如此密集，以至于分开的米尺实际上并不存在。在那之前，它们会被挤压在一起。这一讲还提醒我们，相对论无论如何都不允许真正刚性的物体：如果一个物体是完全刚性的，那么在一侧推一下，就会瞬时把信号传到另一侧，比光还快。}
\end{classroomqa}

\begin{classroomqa}
\audiencequestion{你认为为什么会出现牛顿方程和相对论方程的这种巧合？}\lecturetimestamp{01:49:03}
\lecturerresponse{因为牛顿处理只被应用于一个足够小的区域。只要该区域与整个宇宙的曲率半径相比很小，并且只要那里的局部哈勃速度远低于光速，牛顿力学就是一个很好的近似。这一讲将此比作研究地球上草的生长：人们可以不使用整个地球的曲率来理解局部过程。}
\end{classroomqa}

这就是局部小块的逻辑。牛顿力学首先假设空间曲率对当前问题不重要，其次假设局部没有任何东西以接近光速运动。在宇宙学中，这意味着选择一个足够大以具有均匀性的区域，但仍然足够小，使得曲率效应可以忽略，并且哈勃速度保持非相对论性。

考虑到这一点，牛顿计算并非偶然。它是某个方程的局部近似，而该方程也可以以相对论形式来读。如果把曲率项移到左边，弗里德曼方程就变成
\begin{equation}
\left(\frac{\dot a}{a}\right)^2+\frac{k}{a^2}=\frac{8\pi G}{3}\rho.
\end{equation}
左边是几何的：它涉及标度因子，并且下一次还会涉及空间曲率。右边是能量密度。这正是人们从广义相对论所期望的那种方程。

下一次的主题是曲率本身、爱因斯坦方程，以及那个相对论推导与牛顿推导之间的关系。
